% Synthetic two-column paper for lore-docsrc tests (all content invented).
% Rebuild: pdflatex main.tex && pdflatex main.tex  (then copy main.pdf to ../twocol.pdf)
\documentclass[twocolumn]{article}
\usepackage[margin=0.8in]{geometry}
\usepackage{amsmath}
\pagestyle{empty}
\title{A Toy Model for Testing Asset Extraction}
\author{Ellrin Example and Daidai Sample}
\date{}
\begin{document}
\maketitle
\begin{abstract}
We describe a toy encoder and decoder pipeline that exists only to test the extraction of figures, tables and equations from a two-column paper.
\end{abstract}

\section{Introduction}\label{sec:intro}
Toy pipelines are useful for testing because every element is known in advance. The encoder reads an image and the decoder writes a short report, as shown in Fig.~\ref{fig:arch}. We summarise the results in Table~\ref{tab:res} and define the loss in Eq.~\eqref{eq:loss}.
This paragraph continues with ordinary prose so that the column has enough body text to bound the regions above and below each float, which is what the detector relies on when it searches for the free space next to a caption.

\begin{figure}[t]
\centering
\setlength{\unitlength}{1pt}
\begin{picture}(200,90)
\put(0,20){\framebox(50,50){Encoder}}
\put(50,45){\vector(1,0){40}}
\put(90,20){\framebox(60,50){Memory}}
\put(150,45){\vector(1,0){30}}
\put(180,30){\line(0,1){30}}
\end{picture}
\caption{The toy architecture with an encoder and a memory block.}
\label{fig:arch}
\end{figure}

\section{Method}\label{sec:method}
The encoder maps an input $x$ to features $h$ with a single linear layer, and the decoder reads the features together with a small memory bank $M$ that stores report fragments for later use.
\begin{equation}
h = W x + b
\label{eq:enc}
\end{equation}
The memory is updated with a running average of the features seen during training, which keeps the toy model stable.
\begin{equation}
\mathcal{L} = \sum_{i=1}^{N} \left( y_i - \hat{y}_i \right)^2
\label{eq:loss}
\end{equation}
where $y_i$ is the target and $\hat{y}_i$ is the prediction. Training uses plain gradient descent with a fixed learning rate and no tricks at all, so that the text stays simple for the parser.

\section{Results}\label{sec:results}
Table~\ref{tab:res} lists the scores of the toy model and a baseline on two invented datasets. The toy model is better on both, which is not surprising since the numbers were chosen by hand for this test document.

\begin{table}[h]
\caption{Scores on the invented datasets.}
\label{tab:res}
\centering
\begin{tabular}{lcc}
\hline
Method & Set A & Set B \\
\hline
Baseline & 0.61 & 0.58 \\
Toy model & 0.72 & 0.69 \\
\hline
\end{tabular}
\end{table}

More prose follows the table so that its region is bounded below by body text. The discussion is short because the document has no real findings to discuss beyond the behaviour of the extractor itself.

\begin{thebibliography}{9}
\bibitem{one} E. Example, ``A first toy reference,'' \emph{Journal of Tests}, 2020.
\bibitem{two} D. Sample, ``A second toy reference,'' in \emph{Proc. Imaginary Conf.}, 2021.
\end{thebibliography}
\end{document}
